\section{Introduction}

\begin{frame}{Why verify probabilistic programs?}
  \begin{itemize}
    \item Randomised code is everywhere: cryptography, differential privacy,
          randomised algorithms, probabilistic ML, scientific simulation.
    \item Probabilistic programs \emph{evade conventional testing} --- a rare
          (correct) outlier execution and a bug look the same.
    \item We want \alert{foundational verification}: machine-checked proofs of
          correctness resting on nothing but the axioms of mathematics.
  \end{itemize}
\end{frame}

\begin{frame}{Probabilistic separation logics}
  \begin{itemize}
    \item Separation logic has been hugely successful for heap-manipulating and
          concurrent programs.
    \item \alert{``Separation is independence of random variables.''}
          --- the key idea of Lilac.
    \item Lilac supports \emph{continuous} distributions.
  \end{itemize}
  \pause
  \begin{block}{This project}
    The first mechanisation of a PSL with continuous distributions, and the
    first mechanisation of any PSL in Lean.
  \end{block}
\end{frame}

\begin{frame}{Contributions}
  \begin{itemize}
    \item A complete mechanisation of \textbf{Core Lilac} in Lean 4, with full
          soundness proofs.
    \item An \textbf{interactive proof interface} built on the Lean port of Iris.
    \item Worked examples carried out in the Iris proof mode.
  \end{itemize}
\end{frame}
